\documentclass[11pt,a4paper]{article}

\usepackage[T1]{fontenc}
\usepackage[utf8]{inputenc}
\usepackage{amsmath}
\usepackage{booktabs}
\usepackage{array}
\usepackage{geometry}
\usepackage{hyperref}
\usepackage{lmodern}
\usepackage{longtable}
\usepackage{microtype}
\usepackage{parskip}

\geometry{margin=1in}
\newcolumntype{L}[1]{>{\raggedright\arraybackslash}p{#1}}
\hypersetup{
  colorlinks=true,
  linkcolor=black,
  urlcolor=blue,
  pdftitle={From Refining Capacity to Import Dependence: Portugal's Diesel and Gasoline Market, 2005--2024},
  pdfauthor={Diogo Ribeiro}
}

\title{From Refining Capacity to Import Dependence:\\
Portugal's Diesel and Gasoline Market, 2005--2024}
\author{Diogo Ribeiro}
\date{2026-08-19}

\begin{document}

\maketitle

\begin{abstract}
\noindent
Portugal's refining system gained diesel conversion capability when the Sines hydrocracker
entered commercial production in 2013, and lost one of its two refineries when Matosinhos
ceased refining in May 2021. Using monthly JODI physical flows, the Eurostat oil balance,
DGEG national statistics and the European Commission Weekly Oil Bulletin, this report
measures how the product balance and price exposure changed. The 2013 upgrade produced a
clear structural break: diesel refinery output rose and Portugal moved from a net diesel
importer to a net diesel exporter, with the net-import-to-demand ratio breaking at 2013
($F=19.62$, Benjamini--Hochberg adjusted $p=0.003$). After 2021 the direction reverses.
Annual break tests at 2022 detect nothing, but they have only three post-event observations;
a monthly design with month fixed effects separates the two episodes and finds both to be
distinct. Diesel refinery output falls by $105$ kt per month at the Matosinhos boundary
($p=0.004$) and by a further $152$ kt per month during the 2022 energy-stress period
($p<0.001$), while the net-import-to-demand ratio rises by $0.20$ and $0.38$ respectively.
Short-run pass-through from Spanish to Portuguese pre-tax diesel prices increases from
$0.73$ to $1.01$ after the transition ($p<0.001$). Gasoline short-run pass-through does not
move ($p=0.23$), but an error-correction model, which the cointegration diagnostic indicates
for gasoline, shows its adjustment speed roughly tripling from $-0.091$ to $-0.299$ per week
($p=0.004$), shortening the half-life of a price gap from about eight weeks to two. Both fuels
therefore became more tightly linked to Spain, through different channels. 2022 confounds the
post-closure window and this report does not claim a causal effect of the closure itself.
\end{abstract}

\section{Research Question and Hypotheses}

The empirical question is whether the loss of refining capacity increased Portugal's
dependence on imported finished petroleum products and reduced supply resilience, and
whether Portugal--Spain price co-movement changed after the closure.

Four mechanisms are separated rather than assumed to operate as one chain: refining
capability, the production and trade balance, physical import dependence, and price
co-movement. The pre-specified hypotheses are that the 2013 Sines hydrocracker changed
diesel-producing capability and the trade balance; that the 2021 Matosinhos transition
changed import dependence and domestic refinery-output coverage; that 2022 is a distinct
supply-system stress episode rather than clean post-treatment evidence; and that any price
change must be evaluated on pre-tax prices.

\section{Data and Coverage}

Four independent sources are used. Table~\ref{tab:coverage} states what each covers over the
2005--2024 study window.

\begin{table}[htbp]
\centering
\caption{Source coverage over the study window.}
\label{tab:coverage}
\begin{tabular}{L{3.0cm}L{4.4cm}L{5.8cm}}
\toprule
\textbf{Source} & \textbf{Role} & \textbf{Coverage 2005--2024} \\
\midrule
JODI Oil World Database & Monthly physical flows: imports, exports, refinery output, demand & Complete. 160 of 160 annual product-flow cells, every year with 12 observed months \\
Eurostat \texttt{nrg\_cb\_oil} & Annual balance; Portugal and Spain & Complete for both countries, 160 of 160 cells each \\
DGEG trade workbooks & National import/export statistics & \textbf{2019--2024 only.} 2005--2018 is not published in this series \\
DGEG long sales workbook & Domestic market sales & Complete, 1970--2024 \\
EC Weekly Oil Bulletin & Weekly pre-tax and post-tax consumer prices & Complete. 995 Portuguese and 994 Spanish weekly observations per product \\
\bottomrule
\end{tabular}
\end{table}

\subsection{The one material coverage gap}

DGEG publishes annual product trade workbooks only from 2019. For 2005--2018 the trade
cross-check therefore runs against Eurostat rather than against the DGEG primary publication.
This is a weaker form of corroboration than it appears, because Eurostat's Portuguese oil data
is compiled from the DGEG national submission and tracks it to a median difference of $0.00\%$
across the 2019--2024 cells where both are available. Eurostat is best understood as a second
channel for the same national statistics, not as a third independent opinion. Domestic demand
is unaffected: the DGEG long sales workbook covers the whole window, so demand has three
sources throughout.

\subsection{Product definitions}

Diesel is JODI \texttt{GASDIES} and Eurostat \texttt{O4671}, gas oil and diesel oil. Gasoline
is JODI \texttt{GASOLINE} and Eurostat \texttt{O4652}, motor gasoline. Aviation gasoline is
excluded. In the DGEG sales workbook, gasoline is the \emph{Gasolinas} total and diesel is
road gasoil plus coloured and marked gasoil; road gasoil alone sits about ten per cent below
the JODI and Eurostat demand concept. Differing biofuel treatment does not drive the residual
differences between sources: the Eurostat blended-biofuel wedge is $0.0$ kt on exports in
every year of the window and at most $4.5\%$ on imports.

\section{Methods}

For fuel $j$ in year $t$, with imports $M$, exports $X$, domestic demand $D$ and refinery
output $Q^{\text{refinery}}$:

\[
\text{NetImports}_{j,t}=M_{j,t}-X_{j,t},
\qquad
\text{GrossImportDependence}_{j,t}=\frac{M_{j,t}}{D_{j,t}},
\]
\[
\text{NetImportToDemandRatio}_{j,t}=\frac{M_{j,t}-X_{j,t}}{D_{j,t}},
\qquad
\text{RefineryOutputToDemandRatio}_{j,t}=\frac{Q^{\text{refinery}}_{j,t}}{D_{j,t}}.
\]

The refinery-output-to-demand ratio is not self-sufficiency. Portuguese refinery output may be
exported while domestic demand is met partly by imports, and for gasoline this is the normal
state of the system.

Annual structural breaks use Chow tests on a linear trend at pre-specified years, with 2021
excluded as a transition year so it is assigned to neither side, and Benjamini--Hochberg
adjustment across the sixteen tests. Monthly event timing uses a segmented regression on 293
monthly observations with month fixed effects, a calendar-month trend and HAC(3) standard
errors. Each phase-trend interaction is centred on its phase boundary, so a phase coefficient
is the level shift at that boundary against the extrapolated pre-closure trend, and its
companion trend term is the change in slope within the phase. Price model family is selected
from stationarity and cointegration diagnostics rather than assumed, with HAC(8) standard
errors.

\section{Physical Balance Evidence}

\subsection{Long-run description}

Table~\ref{tab:diesel} gives the annual diesel panel. Two regimes are visible. From 2005 to
2012 Portugal is a net diesel importer, with refinery output covering 73--92 per cent of
demand. From 2013 the ratio exceeds one and the net-import-to-demand ratio turns negative:
Portugal manufactures more diesel than it consumes and exports the surplus. From 2021 the
system returns to net import dependence, and by 2023 refinery output covers only 66 per cent
of demand, the lowest value in the series.

\begin{table}[htbp]
\centering
\caption{Portugal diesel annual balance, kt and ratios. Source: \texttt{fuel\_annual\_analytical\_panel.csv}.}
\label{tab:diesel}
\small
\begin{tabular}{rrrrrrrr}
\toprule
Year & Imports & Exports & Demand & Refinery out. & Gross imp.\ dep. & Net imp./dem. & Ref.\ out./dem. \\
\midrule
2005 & 875 & 211 & 5{,}632 & 5{,}043 & 0.16 & 0.12 & 0.90 \\
2006 & 635 & 314 & 5{,}471 & 5{,}055 & 0.12 & 0.06 & 0.92 \\
2007 & 775 & 193 & 5{,}530 & 4{,}745 & 0.14 & 0.11 & 0.86 \\
2008 & 1{,}006 & 164 & 5{,}450 & 4{,}671 & 0.18 & 0.15 & 0.86 \\
2009 & 1{,}475 & 94 & 5{,}537 & 4{,}074 & 0.27 & 0.25 & 0.74 \\
2010 & 1{,}167 & 35 & 5{,}601 & 4{,}273 & 0.21 & 0.20 & 0.76 \\
2011 & 1{,}315 & 120 & 5{,}191 & 3{,}782 & 0.25 & 0.23 & 0.73 \\
2012 & 892 & 349 & 4{,}678 & 4{,}059 & 0.19 & 0.12 & 0.87 \\
\midrule
2013 & 581 & 1{,}692 & 4{,}535 & 5{,}687 & 0.13 & $-$0.24 & 1.25 \\
2014 & 697 & 1{,}295 & 4{,}697 & 5{,}117 & 0.15 & $-$0.13 & 1.09 \\
2015 & 783 & 2{,}286 & 4{,}941 & 6{,}279 & 0.16 & $-$0.30 & 1.27 \\
2016 & 942 & 2{,}076 & 4{,}834 & 5{,}901 & 0.19 & $-$0.23 & 1.22 \\
2017 & 897 & 2{,}082 & 4{,}974 & 5{,}984 & 0.18 & $-$0.24 & 1.20 \\
2018 & 827 & 1{,}140 & 5{,}089 & 5{,}228 & 0.16 & $-$0.06 & 1.03 \\
2019 & 1{,}075 & 899 & 5{,}075 & 4{,}777 & 0.21 & 0.03 & 0.94 \\
2020 & 875 & 1{,}304 & 4{,}403 & 4{,}606 & 0.20 & $-$0.10 & 1.05 \\
\midrule
2021 & 1{,}491 & 1{,}378 & 4{,}728 & 4{,}199 & 0.32 & 0.02 & 0.89 \\
2022 & 1{,}275 & 331 & 4{,}962 & 3{,}941 & 0.26 & 0.19 & 0.79 \\
2023 & 1{,}597 & 285 & 5{,}084 & 3{,}348 & 0.31 & 0.26 & 0.66 \\
2024 & 1{,}222 & 731 & 4{,}868 & 4{,}133 & 0.25 & 0.10 & 0.85 \\
\bottomrule
\end{tabular}
\end{table}

Gasoline behaves differently throughout. Portugal is a net gasoline exporter in every year of
the window, with refinery output covering between 1.4 and 2.6 times domestic demand. This is
the asymmetry the ratio definition warns about: a system can manufacture a large surplus of
one light product while importing another.

\subsection{The 2013 upgrade is a clear structural break}

Table~\ref{tab:breaks} reports the Chow tests. Three diesel outcomes break at 2013 and survive
multiplicity adjustment comfortably. The break is in the direction the engineering implies:
more diesel manufactured, more exported, and a net-import position that turns negative.

\begin{table}[htbp]
\centering
\caption{Chow tests at pre-specified break years, 2021 excluded as a transition year.
Source: \texttt{structural\_break\_tests.csv}.}
\label{tab:breaks}
\begin{tabular}{llrrrr}
\toprule
Product & Outcome & Break & $F$ & $p$ & BH-adjusted $p$ \\
\midrule
diesel & net import / demand & 2013 & 19.62 & 0.0002 & 0.0026 \\
diesel & refinery output & 2013 & 16.68 & 0.0003 & 0.0027 \\
diesel & exports & 2013 & 13.79 & 0.0008 & 0.0041 \\
gasoline & refinery output & 2013 & 5.23 & 0.0232 & 0.0867 \\
diesel & imports & 2013 & 4.95 & 0.0271 & 0.0867 \\
\midrule
\multicolumn{6}{l}{\emph{No 2022 break test reaches significance; the smallest adjusted $p$ is 0.378.}} \\
\bottomrule
\end{tabular}
\end{table}

The absence of a detected 2022 break should not be read as evidence that nothing happened.
Each 2022 test has eight pre-event and three post-event annual observations, which gives very
low power against anything short of an enormous effect. The monthly design below is the
appropriate instrument for that period.

\subsection{Monthly evidence separates the closure from the 2022 shock}

Table~\ref{tab:monthly} reports the segmented monthly model for diesel. Both phases carry
independent, statistically distinguishable effects, which is the central reason a monthly
design was required: annual data cannot separate a May 2021 event from a March 2022 one.

\begin{table}[htbp]
\centering
\caption{Segmented monthly event model, diesel, $n=293$ months, month fixed effects, HAC(3).
A phase coefficient is the level shift at that phase boundary.
Source: \texttt{monthly\_event\_models.csv}.}
\label{tab:monthly}
\begin{tabular}{llrrr}
\toprule
Outcome & Term & Estimate & Std.\ error & $p$ \\
\midrule
Refinery output (kt/month) & Matosinhos transition & $-105.21$ & 36.02 & 0.004 \\
                           & Energy stress 2022    & $-151.77$ & 23.31 & $<0.001$ \\
Imports (kt/month)         & Matosinhos transition & $42.23$   & 23.40 & 0.071 \\
                           & Energy stress 2022    & $28.36$   & 11.22 & 0.012 \\
Net import / demand        & Matosinhos transition & $0.2005$  & 0.0713 & 0.005 \\
                           & \quad phase trend     & $0.0360$  & 0.0153 & 0.019 \\
                           & Energy stress 2022    & $0.3755$  & 0.0554 & $<0.001$ \\
\bottomrule
\end{tabular}
\end{table}

Diesel refinery output falls by roughly $105$ kt per month at the May 2021 boundary and by a
further $152$ kt per month in the 2022 stress period. The net-import-to-demand ratio rises by
$0.20$ at the closure and by $0.38$ during the energy-stress phase, and the closure phase also
carries a significant upward slope of $0.036$ per month. Both episodes move the system in the
same direction, and the 2022 effect is the larger of the two.

\subsection{2022 as a stress episode}

Against a 2013--2019 baseline, 2022 diesel exports of $331$ kt sit $2.44$ standard deviations
below the mean and at the zeroth percentile of the baseline distribution, while diesel imports
of $1{,}275$ kt sit $2.74$ standard deviations above and at the hundredth percentile. Gasoline
exports in the same year are unremarkable ($z=-0.46$). The stress is diesel-specific, which is
consistent with the documented suspension of Russian oil-product and vacuum gas oil purchases
feeding Sines diesel manufacturing, and inconsistent with a general refining shutdown.

\subsection{Descriptive change around the transition}

Comparing the five years before 2021 with the three years after, and treating 2021 itself as a
transition year, diesel gross import dependence rises from $0.190$ to $0.274$, a change of
$8.4$ percentage points. The diesel net-import-to-demand ratio rises from $-0.119$ to $0.183$,
a change of $30.2$ percentage points, and refinery output coverage falls from $1.088$ to
$0.767$. These are descriptive pre/post differences, not effect estimates.

\section{Price Co-Movement}

Model family was selected from diagnostics rather than assumed. For diesel, neither Portuguese
nor Spanish pre-tax price levels reject a unit root ($p=0.124$ and $p=0.242$) and the pair is
not cointegrated ($p=0.240$), so a levels regression would be spurious and the short-run
log-difference model carries the inference. For gasoline the levels are also non-stationary but
the pair is cointegrated ($p=0.000026$), so an error-correction model is indicated and
estimated below. No ECM is fitted for diesel: with cointegration rejected, the disequilibrium
term would not be a valid long-run residual.

\begin{table}[htbp]
\centering
\caption{Short-run pass-through, $\Delta \log$ pre-tax price, HAC(8), $n=1{,}077$ weeks.
Sources: \texttt{price\_short\_run\_models.csv}, \texttt{price\_model\_choice.csv}.}
\label{tab:price}
\begin{tabular}{llrrr}
\toprule
Product & Term & Estimate & Std.\ error & $p$ \\
\midrule
diesel   & $\Delta \log$ ES               & 0.7348 & 0.0341 & $<0.001$ \\
         & $\Delta \log$ ES $\times$ post & 0.2719 & 0.0541 & $<0.001$ \\
gasoline & $\Delta \log$ ES               & 0.7997 & 0.0593 & $<0.001$ \\
         & $\Delta \log$ ES $\times$ post & 0.0352 & 0.1375 & 0.798 \\
\bottomrule
\end{tabular}
\end{table}

Weekly pass-through from Spanish to Portuguese pre-tax diesel prices rises from $0.73$ before
the May 2021 transition to approximately $1.01$ after it, and the increase is precisely
estimated. The gasoline short-run interaction is small and statistically indistinguishable from
no change.

That is not the whole picture for gasoline. Table~\ref{tab:ecm} reports the error-correction
model the cointegration diagnostic indicates. The long-run relationship is close to unit
elastic ($\log \text{PT} = 0.147 + 0.974 \log \text{ES}$), and the speed at which Portuguese
gasoline prices close a gap against that relationship roughly triples after the transition,
from $-0.091$ to $-0.299$ per week. The implied half-life of a deviation falls from about eight
weeks to a little over two. Gasoline linkage to Spain therefore did tighten; it tightened in the
long-run adjustment channel rather than in contemporaneous pass-through, which a
difference-only model cannot observe by construction.

\begin{table}[htbp]
\centering
\caption{Gasoline error-correction model, two-step Engle--Granger, HAC(8), $n=1{,}077$ weeks.
Long-run relation $\log \text{PT} = 0.147 + 0.974 \log \text{ES}$.
Source: \texttt{price\_ecm\_models.csv}.}
\label{tab:ecm}
\begin{tabular}{lrrr}
\toprule
Term & Estimate & Std.\ error & $p$ \\
\midrule
Disequilibrium (lagged) & $-0.0907$ & 0.0354 & 0.010 \\
\quad $\times$ post & $-0.2080$ & 0.0720 & 0.004 \\
$\Delta \log$ ES & 0.7889 & 0.0720 & $<0.001$ \\
\quad $\times$ post & 0.1970 & 0.1648 & 0.232 \\
\bottomrule
\end{tabular}
\end{table}

The disequilibrium term is a generated regressor, so the second-stage standard errors are
conditional on the first-stage cointegrating estimate; HAC covariance mitigates but does not
remove this.

\section{Portugal--Spain Comparison}

The mean Portugal-minus-Spain pre-tax diesel spread moves from $+4.12$ to $-26.66$ EUR per
1000 litres across the transition, a change of $-30.78$; the gasoline spread moves by $-18.80$.
These are descriptive. The diesel spread is not stationary (ADF $p=0.605$), and the gasoline
spread does not reject a unit root at five per cent either ($p=0.066$), so neither spread can be
treated as mean-reverting around a stable equilibrium, and spread levels should not be read as a
disequilibrium measure.

Spain is used here as a comparison series, not as a causal control. No pre-trend or common-shock
diagnostics have been estimated that would license a difference-in-differences interpretation,
and both countries were exposed to the same European energy shock in 2022.

\section{Source Reconciliation}

JODI and DGEG are the only genuinely independent pair. Across the 24 overlapping 2019--2024
trade cells, $66.7\%$ agree outright within tolerance. Eight cells diverge and each is recorded
with the series that Eurostat identifies as the outlier: six where JODI is out of line and two
where the DGEG workbook is. Over the full window, JODI and Eurostat agree on $89\%$ of 80 cells,
with only three of the nine divergences occurring before 2019.

A row is treated as divergent only when it breaches both the absolute and the percentage
tolerance. Applying them as alternatives makes a 25 kt floor bind on series of roughly 1{,}600
kt, an effective tolerance of $1.6\%$, and flags agreement as close as $2.1\%$ as failure.

\section{Limitations}

\begin{enumerate}
\item \textbf{2022 confounds the post-closure window.} The Matosinhos transition and the
European energy shock are fourteen months apart. The monthly model separates them
statistically, but COVID-recovery demand, Sines maintenance and the Russian VGO disruption are
competing explanations that are not separately identified.
\item \textbf{No causal design.} No control group, instrument or counterfactual is used. Every
post-2021 result in this report is an association measured around a documented event.
\item \textbf{DGEG trade covers only 2019--2024.} The 2013 break is corroborated against
Eurostat, which is compiled from the DGEG submission rather than being independent of it.
\item \textbf{Low annual power after 2021.} Three post-event annual observations cannot detect
anything but a very large break, so the null 2022 annual results are uninformative rather than
negative.
\item \textbf{Gasoline prices need an ECM.} The cointegration diagnostic indicates an
error-correction specification that has not been estimated, so the reported gasoline short-run
coefficients are incomplete.
\item \textbf{JODI assessment codes unverified.} All Portuguese observations carry assessment
code 1, whose semantics have not been confirmed against JODI documentation.
\item \textbf{Retail, not wholesale, prices.} The Weekly Oil Bulletin reports consumer prices.
Pass-through results describe retail co-movement, not refinery-gate or international wholesale
transmission.
\end{enumerate}

\section{Conclusions}

The 2013 Sines hydrocracker produced an unambiguous structural change. Diesel refinery output
and exports broke upward at 2013 and Portugal moved from net diesel importer to net diesel
exporter, a result that survives multiplicity adjustment across sixteen tests.

After May 2021 the direction reverses, and by 2023 domestic diesel manufacturing covers only
two thirds of demand, the lowest coverage in twenty years. A monthly design separates the
closure from the 2022 energy shock and finds both to be independently associated with lower
refinery output and higher net import dependence, the 2022 effect being the larger.

Short-run pass-through from Spanish to Portuguese pre-tax diesel prices rises from $0.73$ to
approximately $1.01$ across the transition. Gasoline contemporaneous pass-through is unchanged,
but its error-correction speed roughly triples, so both fuels became more tightly linked to
Spanish prices through different channels. The initial reading, that only diesel changed, was an
artefact of applying a difference-only model to a cointegrated pair: that specification cannot
see the long-run adjustment channel where the gasoline change actually appears.

This is an association, not a demonstrated causal effect. The closure and the largest European
energy shock in decades fall within fourteen months of one another, and this design cannot
separate their contributions to the post-2021 level. What the evidence supports is that
Portugal's diesel system moved from surplus to deficit over the study window, that the movement
is concentrated around two documented events, and that Portuguese diesel prices now track
Spanish diesel prices considerably more closely than before.

\section*{Claim--Evidence Matrix}

\small
\begin{longtable}{L{4.3cm}L{3.6cm}L{2.3cm}L{4.2cm}}
\toprule
\textbf{Claim} & \textbf{Evidence file} & \textbf{Evidence level} & \textbf{Caveat} \\
\midrule
\endhead
Sines hydrocracker entered commercial production in 2013 & \texttt{refinery\_events.csv} & documented event & Corporate disclosure, not outcome data \\
Matosinhos refining ceased May 2021 & \texttt{refinery\_events.csv} & documented event & 2021 treated as a transition year \\
Diesel net-import ratio breaks at 2013 & \texttt{structural\_break\_tests.csv} & structural break & Diagnostic, not a causal estimator \\
Portugal was a net diesel exporter 2013--2020 & \texttt{fuel\_annual\_analytical\_panel.csv} & descriptive & Ratio is not self-sufficiency \\
Refinery output falls 105 kt/month at the closure boundary & \texttt{monthly\_event\_models.csv} & association & Quote with its phase-trend term \\
2022 stress adds a further $-152$ kt/month & \texttt{monthly\_event\_models.csv} & association & Confounded with demand recovery \\
2022 diesel exports at the 0th baseline percentile & \texttt{stress\_2022\_metrics.csv} & descriptive & Baseline window chosen ex ante \\
Diesel pass-through rises 0.73 to 1.01 & \texttt{price\_short\_run\_models.csv} & association & Retail, not wholesale, prices \\
Gasoline pass-through unchanged & \texttt{price\_short\_run\_models.csv} & association & ECM indicated but not estimated \\
PT--ES diesel spread moves $-30.8$ EUR/1000L & \texttt{pt\_es\_price\_spread\_summary.csv} & descriptive & Spread is non-stationary \\
Sources agree on 89\% of full-window trade cells & \texttt{jodi\_dgeg\_trade\_reconciliation.csv} & descriptive & Eurostat is DGEG-derived \\
No causal effect of the closure is claimed & --- & --- & 2022 is not separable by design \\
\bottomrule
\end{longtable}

\end{document}
